\documentclass[11pt]{article}
\usepackage[margin=1in]{geometry}
\usepackage{hyperref}
\usepackage{listings}
\usepackage{xcolor}
\usepackage{enumitem}

\lstset{
    basicstyle=\ttfamily\small,
    backgroundcolor=\color{gray!10},
    frame=single,
    framerule=0.5pt,
    breaklines=true,
    showstringspaces=false,
    language=bash,
    morekeywords={brew,cmake,ninja,git,cd,mkdir}
}

\title{Building femm-qt on macOS}
\author{D.\ Meeker}
\date{\today}

\begin{document}
\maketitle

\section{Overview}

femm-qt is a cross-platform port of FEMM (Finite Element Method Magnetics)
using Qt~6 for the GUI and CMake/Ninja for the build system.  This document
describes how to build it from source on macOS (Intel or Apple Silicon).

The build produces six executables:
\begin{itemize}[nosep]
  \item \texttt{femm\_gui} --- the main GUI application
  \item \texttt{femm\_headless} --- headless IPC server (no GUI)
  \item \texttt{fkn} --- magnetics solver
  \item \texttt{belasolv} --- electrostatics solver
  \item \texttt{hsolv} --- heat flow solver
  \item \texttt{csolv} --- current flow solver
\end{itemize}

\section{Prerequisites}

\subsection{Xcode Command Line Tools}

Install the Apple compiler toolchain (Clang, make, git, etc.) from Terminal:
\begin{lstlisting}
xcode-select --install
\end{lstlisting}
Follow the prompts to complete the installation.  This provides
\texttt{clang++}, \texttt{git}, and standard system headers.

\subsection{Homebrew}

If not already installed, install the Homebrew package manager:
\begin{lstlisting}
/bin/bash -c "$(curl -fsSL \
    https://raw.githubusercontent.com/Homebrew/install/HEAD/install.sh)"
\end{lstlisting}
Follow the post-install instructions to add Homebrew to your \texttt{PATH}.

\subsection{CMake and Ninja}

\begin{lstlisting}
brew install cmake ninja
\end{lstlisting}

\subsection{Qt 6}

Install Qt~6 via Homebrew:
\begin{lstlisting}
brew install qt@6
\end{lstlisting}
Homebrew installs Qt to a prefix such as \texttt{/opt/homebrew/opt/qt@6}
(Apple Silicon) or \texttt{/usr/local/opt/qt@6} (Intel).  Note this path;
it is needed for the CMake configure step.

To find your Qt prefix:
\begin{lstlisting}
brew --prefix qt@6
\end{lstlisting}

\section{Getting the Source}

If you have a source zip, unzip it:
\begin{lstlisting}
unzip femm-qt-src.zip
cd femm-qt
\end{lstlisting}

Or, if the source is hosted on GitHub, clone the repository:
\begin{lstlisting}
git clone <repository-url> femm-qt
cd femm-qt
\end{lstlisting}

\section{Building}

From the \texttt{femm-qt} directory, run the build script:
\begin{lstlisting}
./macos_build.sh
\end{lstlisting}

This checks that the build tools are present, locates Qt~6 via
Homebrew, configures the project with CMake, and builds it with Ninja.
The GUI is produced as a macOS application bundle
(\texttt{build/bin/femm\_gui.app}), and the script then assembles
everything else into it: the solvers, \texttt{femm\_headless}, the Lua~4
interpreter, the Python scripting assets, the examples, the language
bindings, the manual and the material libraries.  \texttt{macdeployqt}
bundles the Qt frameworks.

The bundle is the \emph{only} place anything runs from on macOS.  The
loose copies CMake leaves in \texttt{build/bin} are deleted once they
have been copied in, so that a stray \texttt{build/bin/femm\_headless}
--- which would resolve its solvers and material libraries from the
wrong directory --- cannot be run by accident.  The regression suite and
the language bindings both reach into the bundle.

The material libraries are a special case, because the material editor
writes them and a signed bundle cannot be written to.  Read-only copies
ship in \texttt{Contents/Resources}, and FEMM copies them to
\texttt{\textasciitilde/Library/Application Support/FEMM/} the first time
they are used; that is where your edits go, and they are shared by every
copy of FEMM on the machine.

To wipe the build directory and rebuild from scratch, pass
\texttt{clean}:
\begin{lstlisting}
./macos_build.sh clean
\end{lstlisting}

If the build fails because CMake cannot find Qt~6, see
\textit{Manual configuration} under Troubleshooting.

\section{Running}

Launch the GUI by double-clicking \texttt{femm\_gui.app} in Finder, or
from Terminal:
\begin{lstlisting}
open build/bin/femm_gui.app
\end{lstlisting}

Or the headless server (for scripted/IPC use).  It lives inside the
bundle, beside the GUI:
\begin{lstlisting}
./build/bin/femm_gui.app/Contents/MacOS/femm_headless
\end{lstlisting}

\subsection{Setting FEMM\_DIR}

The language bindings (OctaveFEMM, MathFEMM, pyFEMM, JuliaFEMM) locate
the FEMM executables via the \texttt{FEMM\_DIR} environment variable.
Point it at \texttt{build/bin}:
\begin{lstlisting}
export FEMM_DIR=/path/to/femm-qt/build/bin
\end{lstlisting}
Enter the \texttt{export} line exactly as shown --- a single line with no
spaces around the \texttt{=}.  This one setting serves both a GUI session
(\texttt{openfemm()} in OctaveFEMM) and a headless session
(\texttt{openfemm(1)}): the bindings look inside
\texttt{femm\_gui.app/Contents/MacOS} for both executables, so either
\texttt{build/bin} or the bundle's own \texttt{Contents/MacOS} works and
both reach the same binaries.  For an installed copy there is no
\texttt{build/bin}, so use the bundle:
\begin{lstlisting}
export FEMM_DIR=/Applications/femm_gui.app/Contents/MacOS
\end{lstlisting}
Add it to your \texttt{\textasciitilde/.zshrc} (or
\texttt{\textasciitilde/.bash\_profile}) to make it persistent.

The binding files themselves are inside the bundle too, at
\texttt{Contents/Resources/bindings} --- \texttt{mfiles} for OctaveFEMM,
and \texttt{mathfemm}, \texttt{scifemm}, \texttt{juliafemm},
\texttt{pyfemm} and \texttt{luafemm} beside it.

\section{Setting Up OctaveFEMM}

OctaveFEMM drives FEMM from Octave (or MATLAB) over a local socket, and uses
Java to make that connection.  On a clean macOS machine the setup is:

\begin{enumerate}
\item \textbf{Install Octave} via Homebrew:
\begin{lstlisting}
brew install octave
\end{lstlisting}

\item \textbf{Install a JDK.}  OctaveFEMM's socket bridge needs a Java
runtime, and Homebrew's Octave does not bundle one.  Install the Adoptium
Temurin JDK --- either download it from \url{https://adoptium.net} or, more
simply, install it with Homebrew:
\begin{lstlisting}
brew install --cask temurin
\end{lstlisting}

\item \textbf{Set \texttt{FEMM\_DIR} and \texttt{JAVA\_HOME}} in your
\texttt{\textasciitilde/.zshrc}.  \texttt{FEMM\_DIR} points at
\texttt{build/bin} (see ``Setting FEMM\_DIR'' above); \texttt{JAVA\_HOME}
is best resolved dynamically with \texttt{java\_home}, so it survives JDK
updates:
\begin{lstlisting}
export FEMM_DIR="/path/to/femm-qt/build/bin"
export JAVA_HOME="$(/usr/libexec/java_home)"
\end{lstlisting}
Open a new terminal (or run \texttt{source \textasciitilde/.zshrc})
afterward so the variables take effect.

\item \textbf{Add the m-files to Octave's path.}  Start Octave and run, once:
\begin{lstlisting}
addpath('/path/to/femm-qt/octavefemm');
savepath;
\end{lstlisting}
\texttt{savepath} records the directory in your Octave preferences so you do
not have to re-add it every session.
\end{enumerate}

To verify, run \texttt{openfemm} at the Octave prompt: the FEMM window should
appear and connect.  Use \texttt{openfemm(1)} for a headless session (no
window), which talks to \texttt{femm\_headless} instead.

\section{Packaging a Disk Image}

To wrap the built bundle in a draggable \texttt{.dmg}:
\begin{lstlisting}
packaging/macos_dmg.sh
\end{lstlisting}

This is a separate step from the build, in the same way that running
Inno Setup is separate from compiling on Windows.  The bundle has to
exist to run and to test; a disk image is only needed to hand FEMM to
someone else.  The script refuses to run if \texttt{femm\_gui.app} is
missing, incomplete, or stale --- the last of those catches a
\texttt{ninja} that relinked the GUI without reassembling the bundle
around it.

The image is written to \texttt{packaging/output/} and contains the
application and a symlink to \texttt{/Applications}, so installing is
the usual drag.  It also carries an \texttt{Examples} folder.

That folder is not a duplicate of the copy inside the bundle, and the
distinction matters.  Example scripts open a model and immediately save
it under another name, so that the original is never modified:
\begin{lstlisting}
open("Roters-Ch9Fig6.fem")
mi_saveas("temp.fem")
\end{lstlisting}
Both paths are relative to the script's own directory, which is where
FEMM sets the working directory when it runs a script.  Inside the
bundle that directory is read-only and the save fails, so drag the
\texttt{Examples} folder somewhere of your own and work there.  The
read-only copy inside the bundle is a master, kept so the examples
survive throwing the disk image away.

\subsection{Signing and notarisation}

\texttt{macos\_build.sh} applies an ad-hoc signature, which is enough to
build, run and test on your own machine.  It is \emph{not} enough for a
disk image handed to anyone else: downloaded from the internet, an
ad-hoc signed application is refused by Gatekeeper as coming from an
unidentified developer.

Distributing to others needs an Apple Developer ID, with which the
bundle is signed, the image notarised, and the ticket stapled:
\begin{lstlisting}
codesign --force --options runtime --timestamp \
         --sign "Developer ID Application: ..." \
         build/bin/femm_gui.app
xcrun notarytool submit <image>.dmg \
         --keychain-profile <profile> --wait
xcrun stapler staple <image>.dmg
\end{lstlisting}
Those steps are deliberately not part of \texttt{macos\_dmg.sh}: they
need credentials, a network round trip to Apple, and decisions about the
hardened runtime and entitlements that are best made against real error
messages rather than guessed at in advance.

\section{Troubleshooting}

\subsection{Manual configuration}

The build script wraps the CMake/Ninja steps and assembles the app
bundle.  The order of the bundling steps matters: \texttt{macdeployqt}
must run \emph{before} the solvers and data files are copied in (it scans
\texttt{MacOS/} and aborts on the \texttt{.dat} files otherwise), and the
bundle must be re-signed afterward.  The material libraries are the one
thing FEMM writes beside itself, so they are seeded read-only into
Contents/Resources/ and copied into the user's Application Support
folder on first use; a signed bundle is never written to.
To run the steps by hand:
\begin{lstlisting}
mkdir build && cd build
cmake .. -G Ninja \
    -DCMAKE_BUILD_TYPE=Release \
    -DCMAKE_PREFIX_PATH="$(brew --prefix qt@6)"
ninja
APP=bin/femm_gui.app/Contents/MacOS
RES=bin/femm_gui.app/Contents/Resources
# 1. Bundle Qt frameworks FIRST, while the bundle holds only femm_gui.
$(brew --prefix qt@6)/bin/macdeployqt bin/femm_gui.app
# 2. Copy the solvers and the headless engine in afterward:
for s in fkn belasolv hsolv csolv femm_headless; do
    cp bin/$s "$APP/"
done
# 3. Material libraries: read-only SEEDS in Resources/, not MacOS/.
#    FEMM copies them to ~/Library/Application Support/FEMM/ on first
#    use and reads and writes them there, so nothing is ever written
#    inside the bundle.
mkdir -p "$RES" && cp ../matlibs/*.dat "$RES/"
# 4. Re-seal the bundle with a deep ad-hoc signature:
codesign --force --deep --sign - bin/femm_gui.app
\end{lstlisting}
The \texttt{macos\_build.sh} script is the authoritative, commented version
of this procedure.

\subsection{Qt not found}

If CMake reports that it cannot find Qt6, verify the prefix path:
\begin{lstlisting}
brew --prefix qt@6
\end{lstlisting}
Pass the result to \texttt{-DCMAKE\_PREFIX\_PATH}.

\subsection{Permission denied on first run}

macOS Gatekeeper may block unsigned applications.  If you see a
``cannot be opened because the developer cannot be verified'' dialog,
go to System Settings $\rightarrow$ Privacy \& Security and click
``Allow Anyway,'' or remove the quarantine attribute:
\begin{lstlisting}
xattr -dr com.apple.quarantine build/bin/femm_gui.app
\end{lstlisting}

\end{document}
